% ID: FIS-TER-MAC-001
% Titolo: Temperature delle sorgenti di una macchina termica
% Disciplina: Fisica
% Area: Termologia
% Argomento: Macchine termiche
% Sottoargomento: Rendimento e macchina di Carnot
% Classe: Quarta liceo scientifico
% Difficolta: 4
% Tipo: problema_numerico
% Risultato: soluzione da aggiungere
% Tempo_stimato: 10 min
% Competenze: modellizzazione, uso del rendimento di Carnot, calcolo algebrico
% Prerequisiti: macchine termiche, rendimento, ciclo di Carnot, temperatura assoluta
% Tag: termologia, macchine_termiche, carnot, rendimento, temperature
% Fonte: verifica_fisica_gas_perfetti_anonimizzata
% Autore: Riccardo Carli
% Licenza: CC-BY-SA-4.0

\begin{esercizio}
Una macchina termica ha rendimento \(8{,}4\%\), pari al \(72\%\) del rendimento
di una macchina di Carnot che opera tra le stesse temperature. La differenza
tra le temperature delle due sorgenti e \(35\,^\circ\mathrm{C}\).

Calcola le temperature delle due sorgenti.
\end{esercizio}

\begin{soluzione}
% Soluzione da aggiungere.
\end{soluzione}

